\chapter{Species and analytic functors}
\label{ch:species}

% Part V, added 2026-09-25 from ~/Downloads/chygraph_master_equation/draft.tex,
% Sec. 7. Nothing is computed here. The two things this chapter shows by
% calculation -- that the excess functions are derivatives, and that the
% Jacobian at threshold is the second derivative -- are identities between
% generating functions the book already carries; the rest is where the same
% objects live in combinatorics, and what would follow if the identification
% were carried through.

Equation~\eqref{eq:act} is a generating function acting on a law through a
family of symmetric multilinear maps. Combinatorics has a name for a
generating function whose coefficients are families of structures acted on
by the symmetric groups: it is the \emph{analytic functor} of \citet{joyal1986},
the functor a species of structures induces, and Eq.~\eqref{eq:act} is that
functor evaluated on a probability space that carries a symmetric multilinear
structure. Once the identification is made, three things the book found by
calculation acquire a reason.

The first is that the excess bracket is a derivative. Species have a
derivative, it is a structure with one extra distinguished element, and
``arrived along an inclusion'' is that element. The second is that the
action on laws is a species composed with the functor that sends a space to
its probability measures, so that the two steps of Chapter~\ref{ch:statmech}
--- a convolution and a pushforward --- are one construction with two
choices of the multilinear map. The third is that the pair is a
Dyson--Schwinger equation in the sense of renormalisation, whose solutions
are sums over rooted trees, which is what a cavity recursion is when written
out. The chapter takes the three in that order, and then says what each
gives the book and what the book gives back.

\section{The connection}
\label{sec:sp-connection}
\index{species|textbf}
\index{analytic functor|textbf}

A \emph{species} $F$ \citep{joyal1981} assigns to every finite set $U$ of
labels a set $F[U]$ of structures on those labels --- the graphs on $U$, the
rooted trees on $U$, the subsets of $U$ --- in a way that respects
relabelling: a bijection $U\to V$ carries $F[U]$ to $F[V]$. Its exponential
generating function is $\sum_{n}|F[n]|\,x^{n}/n!$, and the point of the
theory is that the operations one performs on generating functions --- sum,
product, composition, derivative --- are shadows of operations on the
structures themselves, so that an identity between series is a bijection
between structures. \citet{bergeron1998} is the standard account, and its
title, \emph{Combinatorial Species and Tree-like Structures}, is the subject
of this book in four words.

The generating functions of Sec.~\ref{sec:ensembles} are species in this
sense with one adjustment. $\Phi^{l}_{k}(x)=\sum_{n}p^{lk}_{n}x^{n}$ has for
its coefficients probabilities rather than counts of structures. That is a
\emph{weighted} species \citep{bergeron1998}: the species of finite
sets --- a layer-$l$ complex's set of layer-$k$ containers --- with each set
of $n$ elements carrying the weight $n!\,p^{lk}_{n}$, so that the weighted
exponential series is $\Phi^{l}_{k}$. Nothing in the calculus of species
depends on the weights being counts, and every statement below about
derivatives and compositions holds for weighted species verbatim.

Joyal's analytic functor sends a set $X$ to the set of $F$-structures whose
labels are elements of $X$,
\begin{equation}
  F(X)=\sum_{n}F[n]\times_{S_{n}}X^{n},
  \label{eq:sp-analytic}
\end{equation}
the quotient by $S_{n}$ forgetting the order in which the $n$ labels were
listed. Set this against Eq.~\eqref{eq:act},
$F\act{T}P=\sum_{n}f_{n}(T_{n})\push P^{\otimes n}$. The set $X$ has become
a probability measure $P$ on the message space; the power $X^{n}$ has become
the product measure $P^{\otimes n}$; the quotient by $S_{n}$ is automatic
because the maps $T_{n}$ are symmetric; and the pushforward through $T_{n}$
lands the result back in the message space, where the next application can
act on it. The action of a generating function on a law that
Sec.~\ref{sec:substitution} introduced by fiat, and Chapter~\ref{ch:chyequation}
generalised to any multilinear map, is therefore the analytic functor of the
chy-degree species, evaluated not on a set but on a law.

Three consequences follow, and Secs.~\ref{sec:sp-holes}, \ref{sec:sp-giry}
and~\ref{sec:sp-hopf} take them one at a time: the excess bracket is a
derivative, the action on laws is a species composed with a monad, and the
pair is a Dyson--Schwinger equation.

\section{One derivative, two holes}
\label{sec:sp-holes}
\index{species!derivative of|textbf}
\index{excess generating function}
\index{multisort species|textbf}

The derivative of a species is defined on structures, not on series:
\begin{equation}
  F'[U]=F[U\cup\{\ast\}],
  \label{eq:sp-derivative}
\end{equation}
an $F$-structure on the labels $U$ together with one extra, distinguished
element $\ast$ that is not a label --- the \emph{hole}. Its generating
function is $dF/dx$: a structure of size $n+1$ with one of its $n+1$ elements
singled out and removed from the count contributes $(n+1)x^{n}$. That is what
``arrived along an inclusion'' means. A complex reached along an inclusion is
a complex with one slot already spoken for, and the excess functions of
Sec.~\ref{sec:ensembles} are the derivatives, normalised to be probability
generating functions:
\begin{equation}
  \bar\Phi(x)=\frac{\Phi'(x)}{\Phi'(1)},\qquad
  \bar G(x)=\frac{G'(x)}{G'(1)} .
  \label{eq:sp-excessderiv}
\end{equation}
The book has been using Eq.~\eqref{eq:sp-excessderiv} since
Chapter~\ref{ch:chygraphs} without calling it a derivative, and the moments
of Eq.~\eqref{eq:moments} say so. The first moment of the derivative species
is the second derivative of the original,
\[
  \bar\Phi'(1)=\frac{\Phi''(1)}{\Phi'(1)}
              =\frac{\ave{\kappa(\kappa-1)}}{\ave{\kappa}}=\ave{\kbar},
\]
which is Eq.~\eqref{eq:excess} read from right to left. The size-biased
average that Sec.~\ref{sec:ensembles} explained in words --- an author
picked at random has $\ave{\kappa}_{01}$ papers, an author reached by walking
in along a paper has $\ave{\kbar}_{01}+1$ of them, because the paper walked in
along is the hole and is not among those then found --- is the statement
that pointing a structure and then counting its elements is differentiating
its generating function. The publications chygraph of Sec.~\ref{sec:paper} is
a species with two sorts of label, authors and papers, and walking in along a
paper is a derivative with respect to the paper sort.

\paragraph{The excess bracket is a partial derivative.} That last remark is
the general case. Layers are sorts, and a layered ensemble is a
\emph{multisort} species in the sense of \citet{bergeron1998}: a
layer-$l$ complex's containers are labels of $L$ sorts, one per layer they
can lie in, and the joint generating function $\Phi^{l}(x_{1},\dots,x_{L})$
of their numbers by sort is a series in $L$ variables. A complex reached
through a layer-$m$ container has its hole among the sort-$m$ labels, so its
container counts are generated by
\begin{equation}
  \frac{\del_{m}\Phi^{l}(x_{1},\dots,x_{L})}
       {\del_{m}\Phi^{l}(1,\dots,1)},
  \qquad \del_{m}=\frac{\del}{\del x_{m}},
  \label{eq:sp-multisort}
\end{equation}
the partial derivative with respect to the sort arrived through. When the
counts are independent across sorts, $\Phi^{l}=\prod_{k}\Phi^{l}_{k}(x_{k})$,
the partial derivative differentiates the $k=m$ factor and leaves the others
alone, and Eq.~\eqref{eq:sp-multisort} is
$\prod_{k}[\bar\Phi^{l}_{k}]_{k=m}(x_{k})$: excess on the layer arrived
through, ordinary on every other, which is the bracket of
Fig.~\ref{fig:themap} and of every equation since. The sentence ``discount
the inclusion you came in by, and only from routes back into that layer'' is
the sentence ``differentiate with respect to that sort''. When the counts
are not independent, which is the situation Sec.~\ref{sec:joint} constructs
--- three joint distributions with the same marginals and thresholds sixty
per cent apart --- Eq.~\eqref{eq:sp-multisort} is still right and the product
form is not, and the multisort derivative is what the joint calculation of
that section uses without the name.

\paragraph{Two kinds of hole.} What is new relative to species as they stand
is that a complex has two kinds of slot (Fig.~\ref{fig:sp-holes}). A
complex reached from a member has its hole among the member slots, which is
$\bar G$; a complex reached from a container has its hole among the
container slots, which is $\bar\Phi$. An ordinary species has one kind of
slot per structure and one derivative. A species of nested structures has
up-ports and down-ports, and a derivative in each --- a \emph{directed}
derivative --- and the two excess families of Sec.~\ref{sec:ensembles} are
its two components. That is why the percolation map,
Eqs.~\eqref{eq:Qminus}--\eqref{eq:Qplus}, has two lines where Newman's has
one: each line carries the derivative in one direction and the ordinary
function in the other, and Table~\ref{tab:ce-termbyterm} is the record of
which is which.

\begin{figure}[t]
\centering
\begin{adjustbox}{width=0.86\linewidth}
\begin{tikzpicture}[x=1cm,y=1cm]
% ---- left: hole among the members (arrived from a member): \bar G ----
\begin{scope}
  \node[cpx] (A) at (1.0,1.2) {};
  \node[cpx] (C1) at (0.4,2.4) {};
  \node[cpx] (C2) at (1.6,2.4) {};
  \node[atm] (m1) at (0.0,0) {};
  \node[atm] (m2) at (1.0,0) {};
  \node[circle,draw=black!60,dashed,fill=white,minimum size=2.2mm,inner sep=0pt]
        (m3) at (2.0,0) {};
  \draw[inc] (A)--(C1) (A)--(C2) (A)--(m1) (A)--(m2);
  \draw[incd] (A)--(m3);
  \node[lb,right=2pt of A] {$a$};
  \node[lb,below=1pt of m3] {$\ast$};
  \node[tb,anchor=north,align=center] at (1.0,-0.55)
    {arrived from a member:\\ the hole is a member slot, $\bar G$};
\end{scope}
% ---- right: hole among the containers (arrived from a container): \bar\Phi ----
\begin{scope}[xshift=5.0cm]
  \node[cpx] (A) at (1.0,1.2) {};
  \node[cpx] (C1) at (0.4,2.4) {};
  \node[rectangle,draw=black!60,dashed,fill=white,minimum size=2.9mm,inner sep=0pt,line width=.5pt]
        (C2) at (1.6,2.4) {};
  \node[atm] (m1) at (0.0,0) {};
  \node[atm] (m2) at (1.0,0) {};
  \node[atm] (m3) at (2.0,0) {};
  \draw[inc] (A)--(C1) (A)--(m1) (A)--(m2) (A)--(m3);
  \draw[incd] (A)--(C2);
  \node[lb,right=2pt of A] {$a$};
  \node[lb,above=1pt of C2] {$\ast$};
  \node[tb,anchor=north,align=center] at (1.0,-0.55)
    {arrived from a container:\\ the hole is a container slot, $\bar\Phi$};
\end{scope}
\end{tikzpicture}
\end{adjustbox}
\caption{\textbf{Two kinds of hole.} A complex $a$ drawn as an inclusion
diagram, its containers above and its members below. The derivative of a
species is a structure with one distinguished element removed from the
count, the hole $\ast$, shown dashed. An ordinary species has one kind of
element and one derivative. A complex has two: reached from a member (left)
the hole is a member slot and the remaining members are generated by
$\bar G$; reached from a container (right) the hole is a container slot and
the remaining containers by $\bar\Phi$. The two excess families of
Sec.~\ref{sec:ensembles} are the two components of this directed
derivative, and the two lines of the percolation map are the two ways of
taking it.}
\label{fig:sp-holes}
\end{figure}

\section{Lagrange inversion}
\label{sec:sp-lagrange}
\index{Lagrange inversion|textbf}
\index{component size distribution}
\index{critical amplitude}

The practical payoff of the species reading is an inversion formula. A
species equation of the form $A=X\cdot R(A)$ --- Cayley's rooted trees are
$A=X\cdot E(A)$, a root and a set of subtrees --- has its coefficients in
closed form, and the equation for the size of the component reached along an
edge is of that shape. On a graph, with $\bar\Phi$ the excess degree
generating function, the generating function $H(x)$ of the number of nodes
in the component reached along an edge satisfies $H(x)=x\bar\Phi(H(x))$
\citep{newman2001}: the node reached, marked by $x$, and a set of further
components, one per excess edge. Lagrange inversion then gives every
coefficient, not only the singular part:
\begin{equation}
  H(x)=x\,\bar\Phi\bigl(H(x)\bigr)
  \quad\Longrightarrow\quad
  [x^{s}]\,H=\frac{1}{s}\,[y^{s-1}]\,\bar\Phi(y)^{s},
  \label{eq:sp-lagrange}
\end{equation}
the full distribution of the size of the component reached along an edge
\citep{bergeron1998}, and with it the distribution of the size of
the component a random node belongs to, $x\Phi(H(x))$.

On a chygraph the equation is a system, one unknown per directed inclusion
type, and it is the percolation map with a marking variable. Let $H^{ml}_{i}$
be the generating function of the number of complexes, by layer, in the
finite component reached at a layer-$l$ complex from layer $m$ in direction
$i$, with $x_{l}$ marking each layer-$l$ complex counted. Then
\begin{equation}
  \begin{aligned}
  H^{ml}_{-}&=x_{l}\prod_{k}\Phi^{l}_{k}\bigl(H^{lk}_{-}\bigr)
              \prod_{k}\bigl[\bar G^{l}_{k}\bigr]_{k=m}\bigl(H^{lk}_{+}\bigr),\\
  H^{ml}_{+}&=x_{l}\prod_{k}\bigl[\bar\Phi^{l}_{k}\bigr]_{k=m}\bigl(H^{lk}_{-}\bigr)
              \prod_{k}G^{l}_{k}\bigl(H^{lk}_{+}\bigr),
  \end{aligned}
  \label{eq:sp-system}
\end{equation}
which at $x\equiv1$ is Eqs.~\eqref{eq:Qminus}--\eqref{eq:Qplus} with $Q$ the
probability that the component is finite. The system is of Lagrange type in
$2L^{2}$ variables, $w_{j}=x_{j}\varphi_{j}(w)$ with one $\varphi$ per
directed inclusion type, and the multivariate inversion of \citet{good1960}
gives its coefficients: for any function $f$ of the unknowns,
\begin{equation}
  [x^{\mathbf n}]\,f(w)
  =[w^{\mathbf n}]\;f(w)\,\prod_{j}\varphi_{j}(w)^{n_{j}}\,
   \det\Bigl(\delta_{ij}-\frac{w_{j}}{\varphi_{i}(w)}
             \frac{\del\varphi_{i}}{\del w_{j}}\Bigr),
  \label{eq:sp-good}
\end{equation}
the determinant being the multivariate replacement for the $1/s$ of
Eq.~\eqref{eq:sp-lagrange}. What Eq.~\eqref{eq:sp-good} returns is the joint
distribution of how many layer-$k$ complexes a finite component contains,
resolved by layer and by the direction in which each was entered.

Chapter~\ref{ch:giant} obtained the leading singular term of that
distribution by hand. Section~\ref{sec:jacobian} linearised the map about the
trivial fixed point, whose Jacobian is the threshold tensor, and
Sec.~\ref{sec:amplitude} carried the expansion to second order to get the
critical amplitude $B$ in $S=B\Lambda+O(\Lambda^{2})$: the moment hierarchy
of that chapter is the statement that each order of the expansion of
Eq.~\eqref{eq:sp-system} about $x=1$ costs one further order of moment data.
The whole distribution is available from the same pair, at any distance from
the threshold, through Eq.~\eqref{eq:sp-good}, and the determinant in it is
the same matrix whose Perron root Sec.~\ref{sec:jacobian} tracks --- the
factor $\delta_{ij}-w_{j}\del_{j}\varphi_{i}/\varphi_{i}$ at $w=1$ is
$I-J$, and its vanishing at threshold is what makes the coefficients of
Eq.~\eqref{eq:sp-good} decay as a power law there. It should be said plainly
that the book does not compute this distribution anywhere: what
Chapter~\ref{ch:giant} reports is $S$ and $B$, and the finite-component
distribution by layer is a calculation Eq.~\eqref{eq:sp-good} makes routine
and the book leaves undone.

\section{The analytic functor and the Giry monad}
\label{sec:sp-giry}
\index{Giry monad|textbf}
\index{stochastic species|textbf}

Return to Eq.~\eqref{eq:sp-analytic}. The Giry monad is the functor that
sends a measurable space $M$ to the space $\mathcal{G}(M)$ of probability
measures on $M$, with the point mass $m\mapsto\delta_{m}$ as its unit and
averaging --- a measure on measures collapsed to a measure --- as its
multiplication. The message space $M$ of any row of
Table~\ref{tab:ce-substitutions} is such a space, and the law $P$ that
Chapter~\ref{ch:chyequation} iterates is a point of $\mathcal{G}(M)$. So
$P\mapsto F\act{T}P$ is an endomap of $\mathcal{G}(M)$, obtained by
evaluating the analytic functor of the species $F$ not on $M$ but on the
object the Giry monad returns from it, with the symmetric multilinear maps
$T_{n}$ supplying what the analytic functor needs to land back in $M$. The
chygraph equation, Eqs.~\eqref{eq:ensplus}--\eqref{eq:ensminus}, is a
fixed-point equation for that endomap on $2L^{2}$ copies of
$\mathcal{G}(M)$, one per directed inclusion type; the ensemble map
Eq.~\eqref{eq:ensemblemap} is the same on $L^{2}$ copies with $T=\otimes$;
the population dynamics of Chapters~\ref{ch:hittingset}
to~\ref{ch:satisfiability} is the endomap iterated on a sample.

\emph{Stochastic species} is a name for the composite --- a species whose
structures are assembled from labels drawn from a law, and whose evaluation
is again a law. It does not appear to have been studied as such. Random
generation of species structures exists, Boltzmann samplers being the
standard instance, but that samples a fixed species with a size parameter
rather than composing a species with a law and asking for the fixed points
of the composite, which is what density evolution does.

\paragraph{Why it might matter.} Species composition has a calculus ---
Fa\`a di Bruno for the composition of series, its plethystic form for the
composition of species --- and the derivative of an analytic functor is the
analytic functor of the derivative species. If that lifts through the Giry
monad, then linearising the ensemble map about a point mass produces the
derivative species automatically, and Sec.~\ref{sec:sp-holes} has just
shown that the derivative species is the excess function.

At the level of generating functions the lift is a two-line calculation
rather than a possibility, and it is worth writing out because it is the
structural reason for something the book found by computation. Perturb the
law about the point mass $\delta_{0}$ at the trivial message,
$P=\delta_{0}+\epsilon\,\nu$, and expand Eq.~\eqref{eq:act} to first order:
by multilinearity and symmetry of $T_{n}$, each of the $n$ factors can carry
the perturbation, so
\[
  F\act{T}(\delta_{0}+\epsilon\nu)
  =F\act{T}\delta_{0}
   +\epsilon\sum_{n}n f_{n}\,(T_{n})\push
     \bigl(\nu\otimes\delta_{0}^{\otimes(n-1)}\bigr)+O(\epsilon^{2}),
\]
and the coefficient $n f_{n}$ is the coefficient of $x^{n-1}$ in $F'$. The
linearisation of the action of $F$ is the action of $F'$, with $T_{n}$
evaluated on one perturbed leg and $n-1$ trivial ones --- which is the
kernel derivative $u'$ of Eq.~\eqref{eq:uprime}, one leg at a time. Applied
to a message that arrived along an inclusion the series is already
$\bar\Phi=\Phi'/\Phi'(1)$, and linearising differentiates it once more:
\[
  \bar\Phi'(1)=\frac{\Phi''(1)}{\Phi'(1)}=\ave{\kbar},
  \qquad
  \bar G'(1)=\frac{G''(1)}{G'(1)}=\ave{\sbar}.
\]
That is why the Jacobian at threshold, Eq.~\eqref{eq:branch} and the tensor
of Sec.~\ref{sec:threshold}, is built from the excess moments $\ave{\kbar}$
and $\ave{\sbar}$ and not from $\ave{\kappa}$ and $\ave{s}$: ``linearise
about the trivial fixed point'' and ``arrive along an inclusion'' are the
same derivative taken twice, and the elementwise reweighting of
Eq.~\eqref{eq:reweight} is the second derivative's leg carrying $u'$. The
functorial statement --- that this is the derivative of an analytic functor
commuting with the Giry monad, so that it holds for any species and any
monad-compatible multilinear structure and not only for the series at hand
--- is the part that remains to be written down, and it is stated here as
a possibility and not a result.

\section{The Hopf side}
\label{sec:sp-hopf}
\index{Connes--Kreimer Hopf algebra|textbf}
\index{Dyson--Schwinger equation|textbf}

The third consequence takes the cavity recursion apart tree by tree. The
Hopf algebra of \citet{connes1998} is the polynomial algebra on rooted
trees, with a coproduct that cuts a tree into a trunk containing the root
and the forest of pruned branches. The grafting operator $B_{+}$ takes a
forest and puts a new root above it, and its defining property is that it is
a Hochschild $1$-cocycle, which is what makes it compatible with the
coproduct \citep{bergbauer2006}: cutting a grafted tree is either cutting
below the new root or cutting inside the branches. A combinatorial
Dyson--Schwinger equation is
\begin{equation}
  X=B_{+}\bigl(f(X)\bigr),
  \label{eq:sp-dse}
\end{equation}
with $f$ a formal power series, and its solution is a formal sum over all
rooted trees with coefficients fixed by $f$: $f$ says how many branches a
vertex may carry and with what weight, and $B_{+}$ closes the recursion by
supplying the vertex.

Newman's equation, the left-hand side of Eq.~\eqref{eq:sp-lagrange}, is
Eq.~\eqref{eq:sp-dse} with $B_{+}$ in the role of the factor $x$ --- the node
reached, placed above the components hanging from it --- and $\bar\Phi$ in
the role of $f$. Its solution is the branching-process tree written out as a
weighted sum over trees, one term per shape of finite component. The cavity
recursion, which assigns a message to each tree by recursing from the
leaves, is then a \emph{character} of the Hopf algebra in the same way that
Feynman rules are: a map from trees to numbers that is multiplicative on
forests, because on a tree the message at the root is a product over the
subtrees below it. The cuts of the coproduct are the decompositions of a
cavity tree into sub-cavity trees, which is the object the recursion is
built on and the object the excess bracket keeps track of --- the cut edge
is the inclusion arrived through.

\citet{foissy2010} classifies systems of such equations, one unknown per
decoration, in the Hopf algebra of trees whose vertices carry decorations.
Layers are decorations, so a layered chygraph with a single direction ---
Part~II's constructions in which nothing is entered from above --- is
covered by that classification as it stands. The nested chygraph needs two
things beyond it.

First, a vertex is reached either from a container or from a member and its
children are of both kinds, so the grafting has to record the direction of
each edge: two grafting operators, one per direction, which are the two
lines of Eq.~\eqref{eq:sp-system}. This may reduce to Foissy's setting by
decorating each vertex with the pair (layer, direction of arrival), which is
the $2L^{2}$ index set of Chapter~\ref{ch:percolation}, and that is why it
should be checked before it is claimed: the classification would then apply
to the whole percolation theory of Part~II, and would say which of its
systems are Hopf subalgebras --- which fixed-point equations close on a
subalgebra of trees --- and which are not.

Second, and this is the genuine gap, $f$ in Eq.~\eqref{eq:sp-dse} is a power
series only when the weight of a tree factorises over the children of each
vertex. For percolation it does, because reachability factorises: that is
the first row of Table~\ref{tab:ce-substitutions}, $\Tint_{l}$ itself a
product, and it is why the character is multiplicative. For a general kernel
the interior operator couples the children --- the emitted field of
Eq.~\eqref{eq:emitted} is not a product of one-field terms --- and the
weight at a vertex is an operation from an operad rather than a coefficient.
Trees whose vertices are decorated by an operad have their own
Connes--Kreimer construction \citep{moerdijk2001,vanderlaan2003}, and that is
where the nested chygraph lives: a tree of operadically decorated vertices,
with a complex as the operation at a vertex. Chapter~\ref{ch:operads} takes
that up.

\paragraph{A remark on renormalisation.} The Birkhoff decomposition that
renormalisation uses on this Hopf algebra separates the character into a
divergent part and a finite part, factor by factor over the coproduct. At
the percolation threshold the sum over finite trees diverges --- the
coefficients of Eq.~\eqref{eq:sp-good} stop being summable when $I-J$ is
singular --- and the giant component is the part that has to be subtracted
to make the finite-component distribution a probability again. Whether the
subtraction is a Birkhoff decomposition in the technical sense, with the
giant component as the counterterm, is not known; it is recorded here as an
analogy and nothing more.

\section{What it gives the book}
\label{sec:sp-gives}

Four things, in decreasing order of how immediately they can be used.

A one-line account of the excess bracket,
Eq.~\eqref{eq:sp-excessderiv}, which the book explains by a figure and a
paragraph in Chapter~\ref{ch:percolation}: excess is the derivative, the
bracket is the partial derivative with respect to the sort arrived through,
and the two excess families are the two components of a directed derivative.
Section~\ref{sec:joint}'s joint distributions, which the product form cannot
hold, are covered by the multisort form Eq.~\eqref{eq:sp-multisort} without
further apparatus.

A name for what Eq.~\eqref{eq:act} is, so that the convolution of
Eq.~\eqref{eq:conv} and the pushforward of the interior sum are seen to be
one construction --- an analytic functor evaluated on a law --- with two
choices of multilinear map, and a derivation of why the Jacobian at threshold
carries the excess moments: the linearisation of an action is the action of
the derivative, and arriving along an inclusion was already one.

The full finite-component distribution by layer, Eq.~\eqref{eq:sp-good},
where Chapter~\ref{ch:giant} has the amplitude. This is the one item on the
list that is a calculation rather than a reading, and it is not done.

And, stated with the caution Sec.~\ref{sec:sp-hopf} gave it, a candidate
meaning for the Birkhoff decomposition at a percolation threshold.

\section{What the field could get}
\label{sec:sp-gets}

Three objects that do not appear to exist in the theory of species, each of
which the book has already used.

\emph{The directed derivative.} A species with two kinds of hole, up-ports
and down-ports, with a derivative in each and a Lagrange inversion in two
hole types, Eq.~\eqref{eq:sp-system} with Eq.~\eqref{eq:sp-good}. Species
theory has multisort species and their partial derivatives, which is the
layer index; what it does not have is a structure whose elements are
themselves structures with elements, so that a hole can be on either side
of an inclusion. The chygraph is the natural example, and its percolation
theory is the calculus of that derivative worked out.

\emph{The stochastic species.} An analytic functor composed with the Giry
monad, with the chygraph equation as its fixed-point equation and density
evolution as its calculus. The questions such an object raises --- when does
the fixed point exist, when is it unique, what is the derivative of the
composite --- are the questions Chapters~\ref{ch:rde}
and~\ref{ch:chyequation} answer in the physics vocabulary: existence by
iteration from the trivial law, uniqueness by the stability of that fixed
point, the derivative by Sec.~\ref{sec:sp-giry}.

\emph{Dyson--Schwinger systems with two grafting operators over operadically
decorated trees.} The Hopf-algebraic setting in which the nested chygraph
with a general kernel would be a system of Eq.~\eqref{eq:sp-dse}'s type.
The physics of Parts~II and~III solves those systems --- thresholds,
amplitudes, free energies --- before the algebra has been written down, and
the algebra, once written, would say which of the book's fixed-point
equations close on a subalgebra and which do not.
